\section{Implementation Details}
\label{app:implementation}

\paragraph{Off-chain Prover.}
The prover is implemented in Rust and uses the \emph{arkworks} toolkit for BN254 pairings. 
The instance is launched with a CLI that generates all three proof kinds (main, receipt, and dispute) described in Sections~\ref{sec:constraints} and~\ref{sec:voter}.
It commits to all polynomials using the \emph{hiding} variant of KZG, where each commitment includes a blinding term in the $\gamma$-scaled SRS basis.
The blinding polynomials are derived from the prover transcript/challenges, while the hiding structure is enabled by the SRS elements.
Additionally, the prover applies \emph{off-domain masking} by adding random multiples of the domain vanishing polynomial to selected witness polynomials.
This preserves on-domain evaluations while randomizing coefficients off-domain.
To reduce proof size and verification cost, each opening batches multiple polynomial evaluations at the \emph{same point}
using a random scalar $\gamma$ into a single KZG witness. On-chain, this yields one pairing check \emph{per opening point},
while still validating many evaluations within each opening.

\paragraph{Calldata encoding.}
Proofs are serialized into a single \texttt{bytes} blob using Ethereum ABI encoding of fields:
\emph{kind}, \emph{commitments array}, \emph{openings array}, \emph{public-input tuple}, and the transcript scalars $(\alpha,\beta,\zeta,r)$.
The CLI computes the 4-byte function selector as \texttt{keccak256(methodSignature)[0..4]} via \texttt{tiny-keccak} when forming transaction calldata for verification.

\paragraph{On-chain verifier.}
The verifier is implemented in Solidity and built with the \texttt{Foundry} toolchain.
It uses the BN254 precompiles (EIP-196/197) for elliptic-curve operations and pairing checks.
There is no SNARK circuit/verifying key; instead, the contract stores KZG verifier parameters derived from the SRS
($g$, $\gamma g$, $h$, and $\beta h$) in \texttt{ZeeperioKzgVk.sol}.
The contract exposes three entrypoints, \texttt{verifyMain}, \texttt{verifyReceipt}, and \texttt{verifyDispute}.
Each calls a shared internal routine \texttt{\_verify(expectedKind, proofCalldata)} which decodes the proof and runs verification.

\paragraph{Custom ABI decoding via \texttt{calldataload}.}
Proofs are passed as a single \texttt{bytes calldata} argument and decoded manually with \texttt{calldataload}.
The helper \texttt{\_readWord(data, offset)} reads 32-byte words from calldata at \texttt{data.offset + offset}.
The proof blob begins with a fixed-size header:
\[
\begin{array}{ll}
\texttt{0x00:}~\texttt{kind} &
\texttt{0x20:}~\texttt{commitmentsOffset} \\
\texttt{0x40:}~\texttt{openingsOffset} &
\texttt{0x60:}~\texttt{publicOffset} \\
\texttt{0x80:}~\alpha &
\texttt{0xA0:}~\beta \\
\texttt{0xC0:}~\zeta &
\texttt{0xE0:}~r
\end{array}
\]
where offsets point into the same \texttt{bytes} blob.
	
\paragraph{Commitments section.}
At \texttt{commitmentsOffset}, the contract decodes a commitment array of length $\ell$, where each element occupies \texttt{0x60} bytes.
This includes a 32-byte word whose low 16 bits store \texttt{uint16 label}, followed by two field elements encoding a BN254 $G_1$ point $(x,y)$.

\paragraph{Openings section.}
	At \texttt{openingsOffset}, the contract decodes the openings array using the ABI \emph{array-of-offsets} layout:
	\[
	\texttt{len} \,\|\, \texttt{ptr}_0 \,\|\, \cdots \,\|\, \texttt{ptr}_{\texttt{len}-1}.
	\]
	Each pointer is a relative offset to an opening tuple. For each opening, the contract reads a fixed prefix:
	\[
	(\texttt{witness.x}, \texttt{witness.y}, \texttt{point}, \gamma, \texttt{labelsOff}, \texttt{valuesOff}, \texttt{blindingsOff}),
	\]
followed by three arrays at the given offsets:
(i) labels encoded as 32-byte words and decoded by truncation to \texttt{uint16[]},
(ii) \texttt{uint256[] values}, and (iii) \texttt{uint256[] blindings}. 
The \texttt{labels} selects which commitments are needed, \texttt{values} contains the claimed evaluations, and \texttt{blindings} contains the hiding terms needed for the hiding-KZG verifier equation.
	
\paragraph{Public inputs section.}
	At \texttt{publicOffset}, the contract decodes:
	\[
	(n,\; C,\; \texttt{sumA},\; \texttt{sumM},\; \texttt{tallyOffset},\; \texttt{disputedCode},\; \texttt{ballotIndex}),
	\]
	where \texttt{tallyOffset} points to the \texttt{uint256[]} tally array. 
	For the main election proof, \texttt{tally.length} must equal $C$.
	For the \texttt{Dispute} proof, the \texttt{disputedCode} and \texttt{ballotIndex} fields are used to bind the proof statement to the specific dispute instance.

\paragraph{Transcript checks and KZG verification.}
After decoding, the verifier recomputes Fiat--Shamir challenges using \texttt{sha256}-based hashing into $\mathbb{F}_r$.
It checks $\alpha$ and $\zeta$ for all proof kinds, and additionally checks $\beta$ for the receipt proof.
It also recomputes $r$ deterministically from the decoded openings (hashing each opening's \texttt{point}, $\gamma$, witness,
and all \texttt{values}/\texttt{blindings}), rejecting if $r$ does not match the prover-supplied field.
Finally, for each opening point, it verifies a KZG pairing equation (one pairing check per opening), where the referenced commitments
are combined using the opening's $\gamma$ and the claimed \texttt{values}/\texttt{blindings}.

\paragraph{Evaluation on Sepolia.}
We deploy to Sepolia because it is free (test ETH can be obtained from public faucets), which makes it practical to repeatedly measure gas.
Deploying on the Ethereum mainnet requires no protocol changes, as the BN254 precompiles specified in EIP-196/197 are available on both networks.

